% !TeX root = ../Tex/main.tex

\newpage
\subsection*{Model 1 - Surprise as Shock}

\textcolor{red}{Citazione Battigalli su perfect and imperfect information}



\vspace*{1em}
The decisive variable in this situation may be the element of surprise (\cite{schellingStrategyConflict1960}).
The two versions of the elephant-human encounter, the one involving the shepherd and the one involving the tourist, are analogous. In the tourist scenario, there is no surprise: the elephant is aware that humans might approach and disturb it, and it knows where to find them.

The situation involving the shepherd, however, contains the element of surprise, which may justify a different behavioral response, namely, the attack. What surprise alone does not explain, however, is why the elephant protected the shepherd after the attack.

To justify this behaviour we introduce empathy, force that shapes the outcome of the game.

However, before jumping to the model, we must ourselves what does empathy have to do with the first case.
We may argue that empathy weakens elephant’s signal. Yet empathy also functions as a force that regulates the intensity of the signal itself.

Is it therefore reasonable to suggest that, in the second case, empathy may also be the mechanism that modulates the variability of the signal?
What is the relationship between these two forces?

Now that the cause has been established, we can turn to the effects.
We observe two primary effects: that of surprise and that of empathy.

\subsection*{Model 2 - Vote Buying}

\subsubsection*{Paper di Grosier e Snyder 1996}

In this section I propose a readaptation of a famous model (\cite{grosecloseBuyingSupermajorities1996}). I intend to expoilt the genious equilibrioum authors found to support the intuition coming from the DAG with another nerrative. In doing so, I believe we are improving the previous model and we are getting closer to the final aim of representing a real world event.

\vspace*{1em}
In their article (\cite{grosecloseBuyingSupermajorities1996}) the authors explain why in the US Congress the number of votes often exceedes the minimum required to make a bill pass. Thay argued that this is due to the interaction between lobbying groups competing for support of different bills. These groups \textit{bribe} Congressmen to convince them. 

\vspace*{1em}
\textcolor{red}{The basic intuition from the paper comes from this example:}

\vspace{1em}
\textit{ "The legislature contains seven members, all indifferent between the two proposals, that is, $v(i) = 0$ for all $i$. Since B will move second and attack the weakest part of A’s coalition, A will offer the same bribe, $a$, to all legislators he bribes. Let $m + 4 \ge 0$ be the size of the coalition A bribes. If $m = 0$, then B can defeat $x$ by paying $a + \varepsilon$ to one member of A’s coalition. Thus, A must set $a \ge W_B$ to win, and A’s total bribes are then $4W_B$. If $m = 1$, then B must pay $a + \varepsilon$ to two members of A’s coalition to defeat $x$. Thus, A must set $a \ge W_B/2$ to win, and his total bribes are then $5W_B/2$. This is less than $4W_B$, so A is better off buying a supermajority of size 5 than a minimal winning coalition. If $m = 2$, then B must pay $a + \varepsilon$ to three members of A’s coalition, so A must set $a \ge W_B/3$ to win, and his total bribes are then $2W_B$. In fact, A minimizes his total payments by bribing all seven legislators. Here, $m = 3$, and A’s total bribes are $7W_B/4$. Clearly, there is an equilibrium where $x$ wins if and only if $W_A \ge 7W_B/4$. If $W_A < 7W_B/4$, then there are only equilibria in which no bribes are paid and $s$ wins."} (\cite{grosecloseBuyingSupermajorities1996})


\subsubsection*{Representation through Vote Buying}
Consider a council composed of several members. These members are not strategic players in the game but rather represent passive agents whose votes can be influenced.  
There are two active players, conceptualized as two groups, each seeking to "buy" the votes of the council members in order to determine the policy agenda.  
Each group possesses a budget, and bills are voted on by the council according to a predefined rule, assumed here to be the simple majority rule.  
This setup corresponds to the classical \textit{vote-buying game} as formalized by (\cite{grosecloseBuyingSupermajorities1996}).

The objective of this analysis is to adapt this framework to the so-called \textit{elephant-shaped} case.  
This adaptation allows the equilibrium structure of the original model to be employed within a cognitive and emotional context.


\begin{itemize}
    \item Let \( K \) denote the total number of \textit{thoughts}.
    \item The two players represent \textit{opposing emotions} that regulate the internal dynamics of the elephant’s mind.  
    "Setting the agenda" refers to determining the behaviour.
    \item Each player’s \textit{budget} corresponds to the \textit{intensity} of its emotion.  
    A higher emotional intensity enables the emotion to "purchase" a larger number of thoughts.
\end{itemize}

Notation:

\begin{itemize}
  \item[] Players: \(N=\{X,Y\}\) or \(N=\{Emotion 1, Emotion 2\}\).
  \item[] Histories: \(H= \big \{\varnothing,\{x\},\{x,y\}\big \}\). Where \(x\) and \(y\) indicate vectors of payments in the standard setting. Emotions use a unit of payment (emotional intensity) to corrupt a thought (or a judgement).
  \item[] Player function: \(p(\varnothing)=X\);\quad \(p(x)=Y\)
  \item[] Actions: \(A(\varnothing)=\{x\}\);\quad \(A(x)=\{y\}\) \(\forall\,x\).
  \item[] Each player (emotion) \textcolor{red}{receives} \(V_i\) if elephant acts accordingly.
  \item[] Each player loses \(\sum_{j=1}^{k} x_{j}\) for each unit of intensity spent.
\end{itemize}

%\vspace{1em}
This adapted framework enables a formal representation of conflict as a strategic interaction between emotions competing for cognitive control.  
Each emotion allocates its limited intensity across thoughts, attempting to secure a majority sufficient to determine the elephant’s mental and behavioral trajectory.  
The resulting equilibrium can be interpreted as a stable emotional configuration in which neither emotion can improve its outcome given the other’s allocation strategy. Let \( \mu \) denote the minimal number of thoughts required to \textit{set the agenda}, that is, to achieve behavioral dominance of one emotion over the other. Then:
\[
\mu=\dfrac{k+1}{2}
\]



